\documentclass[10pt]{article}
\usepackage[margin=0.9in]{geometry}
\usepackage{amsmath,amssymb,amsthm,microtype,hyperref}
\newtheorem{theorem}{Theorem}
\title{Formal Forward Coverage of Two Modulo Nine}
\author{Collatz proof project}
\date{September 5, 2026}
\setlength{\emergencystretch}{3em}
\begin{document}
\raggedright
\maketitle
\begin{abstract}
We formalize Corollary 5.8 of Monks et al.: every positive natural
Collatz orbit visits two modulo nine. Applying the result to every tail
gives arbitrarily late visits and a conditional reduction of convergence
to this progression. This is a Lean formalization of a known theorem;
it does not exclude divergence or nontrivial cycles.
\end{abstract}
Write $T(n)=n/2$ for even $n$ and $T(n)=(3n+1)/2$ for odd $n$.
\begin{theorem}
For every $n>0$ and $K\geq0$, some $t\geq K$ satisfies
$T^t(n)\equiv2\pmod9$.
\end{theorem}
\begin{proof}
Our existing formal arithmetic results give a tail outside multiples of
three and an even value on every natural orbit. On the six remaining
residues the possible transitions, with parity suppressed, are
\[
\begin{array}{c|cccccc}
r&1&2&4&5&7&8\\\hline
T(r)&5,2&1,8&2&7,8&8,2&4,8.
\end{array}
\]
Assume a tail avoids residue two. The finite transition constraints force
its value after three steps to have residue eight: residue four must be excluded since
its successor is two, and the remaining non-eight chain is $1,5,7,8$.
Applying this argument to every further tail shows that all values from
time three onward have residue eight. An even value of residue eight has
successor residue four, a contradiction. This proves a visit exists.
Starting the same argument at $T^K(n)>0$ gives the stated time bound.
\end{proof}
The even-value lemma uses finite parity injectivity: if a natural seed
$n$ had only odd steps, its parity prefix of length $k=n+2$ would agree
with $2^k-1$, forcing equality of their residues modulo $2^k$.
Both lie below $2^k$, while $n<2^k-1$, a contradiction.

\paragraph{Convergence reduction and remaining gap.}
If every positive $n\equiv2\pmod9$ reaches one, then every positive seed
reaches one by composition with the hitting time above. That premise is
unproved. In particular, the result cannot be combined directly with the
least-noncontracting-seed restriction: the property
$2^t\leq3^{j_t(n)}$ for all $t$ is not supplied at the selected residue
visit. Coverage and the existence of noncontracting tails are separate
existential statements; their witnesses have not been shown to coincide.

\paragraph{Formal artifacts.}
\texttt{Collatz.ResidueCoverage} proves
\texttt{exists\_two\_mod\_nine},
\texttt{arbitrarily\_late\_two\_mod\_nine}, and
\texttt{reaches\_one\_of\_two\_mod\_nine}.
The existing \texttt{exists\_even\_value} is now public.
The proofs use exact arithmetic and existing parity results, with no
imported literature axiom or finite seed cutoff.

\begin{thebibliography}{1}
\bibitem{monks} K. Monks, K. G. Monks, K. M. Monks, and M. Monks,
\emph{Strongly sufficient sets and the distribution of arithmetic
sequences in the $3x+1$ graph}, arXiv:1204.3904v2 (2012), Corollary 5.8.
\url{https://arxiv.org/abs/1204.3904}.
\end{thebibliography}
\end{document}
